\section{The exact gamma saddle and the all orders theorem}
\label{sec:saddle}
For a real dimension mark $t$, let
\begin{equation}\label{eq:marked}
 B_N(t)=\sum_{m=1}^N T(N,m)e^{tm},\qquad B_N(0)=b_N.
\end{equation}
The real phase on $1\le x\le N$ is
\begin{equation}\label{eq:phase}
 f_N(x)=\log\Gamma(M(x)+N-x)-\log\Gamma(N-x+1)-\log\Gamma(M(x)).
\end{equation}
It agrees \emph{exactly} with the summands: $e^{f_N(m)}=T(N,m)$. Replacing this phase by only its leading logarithmic approximation would lose an unbounded exponential factor; Section~\ref{sec:elementary} quantifies this issue.

\begin{theorem}[Exact saddle expansion]\label{thm:main}
Fix a finite $T\ge0$ and an integer $K\ge1$. Write $L=\log N$ and
\[
 I_N=[N/(2L),4N/L].
\]
For all sufficiently large $N$, uniformly for real $|t|\le T$, there is exactly one root $r=r_N(t)$ in $I_N$ of
\begin{equation}\label{eq:saddle}
 f_N'(r)+t=0.
\end{equation}
It is the unique global maximizer of $f_N(x)+tx$ on $[1,N]$. Put
\begin{equation}\label{eq:standard}
 A=-f_N''(r),\qquad \sigma=A^{-1/2},\qquad
 \lambda_j=f_N^{(j)}(r)A^{-j/2}\quad(j\ge3).
\end{equation}
Then
\begin{align}
 r&=\frac{2N}{L-2\log L+\log2+2-t+O_T(\log L/L)},\label{eq:rscale}\\
 A&=\frac{2N}{r^2}\bigl(1+O_T(L^{-1})\bigr)\sim\frac{L^2}{2N},\label{eq:Ascale}\\
 \lambda_j&=O_{j,T}(N^{1-j/2}).\label{eq:lamscale}
\end{align}
Define $C_0=1$ and, for $\ell\ge1$,
\begin{equation}\label{eq:contraction}
 C_\ell=\sum_{\substack{k_3,\ldots,k_{2\ell+2}\ge0\\
                 \sum_{j=3}^{2\ell+2}(j-2)k_j=2\ell}}
 \left(\sum_{j=3}^{2\ell+2}jk_j-1\right)!!
 \prod_{j=3}^{2\ell+2}\frac{(\lambda_j/j!)^{k_j}}{k_j!}.
\end{equation}
Here $(-1)!!=1$. This is a finite sum, its double factorial is an even Gaussian moment, and $C_\ell=O_{\ell,T}(N^{-\ell})$. The marked count satisfies
\begin{equation}\label{eq:main}
 B_N(t)=e^{f_N(r)+tr}\sqrt{\frac{2\pi}{A}}
 \left(\sum_{\ell=0}^{K-1}C_\ell+O_{K,T}(N^{-K})\right).
\end{equation}
\end{theorem}
The theorem permits each \emph{fixed} $K$ and bounded real $t$. It asserts neither convergence as $K\to\infty$ nor an error uniform in an order growing with $N$. It also does not assert the absence of smaller stationary points outside $I_N$.

For $f_j=f_N^{(j)}(r)$, the first corrections are
\begin{align}
 C_1&=\frac{f_4}{8A^2}+\frac{5f_3^2}{24A^3},\label{eq:C1}\\
 C_2&=\frac{f_6}{48A^3}+\frac{7f_3f_5}{48A^4}
       +\frac{35f_4^2}{384A^4}
       +\frac{35f_3^2f_4}{64A^5}
       +\frac{385f_3^4}{1152A^6}.\label{eq:C2}
\end{align}
In particular,
\begin{equation}\label{eq:C-limits}
 C_1\sim\frac1{24N},\qquad C_2\sim\frac1{1152N^2}.
\end{equation}
The limits in \eqref{eq:C-limits} cannot replace the exact coefficients in \eqref{eq:main} at its stated power accuracy. Their logarithmic relative corrections are larger than the next inverse-$N$ power. For example, in the bounded computations through $N=1000$, the exact $C_1$ at $t=0$ is still negative.

The rest of the proof is split into global localization and a lattice expansion. This makes explicit why neither an isolated formal saddle calculation nor an unproved replacement of the sum by an integral is enough.
